% Generated draft: requires mathematical and source review.
\documentclass[11pt]{article}
\usepackage[margin=25mm]{geometry}
\usepackage{fontspec}
\setmainfont{Latin Modern Roman}
\usepackage{amsmath,amssymb,mathtools,unicode-math}
\setmathfont{Latin Modern Math}
\usepackage{enumitem,needspace,xurl,hyperref}
\setlength{\parindent}{0pt}
\setlength{\parskip}{5pt}
\setlength{\emergencystretch}{3em}
\newcommand{\N}{\mathbb{N}}
\newcommand{\Z}{\mathbb{Z}}
\newcommand{\Q}{\mathbb{Q}}
\newcommand{\R}{\mathbb{R}}
\newcommand{\C}{\mathbb{C}}
\newcommand{\F}{\mathbb{F}}
\newcommand{\A}{\mathbb{A}}
\newcommand{\PP}{\mathbb{P}}
\newcommand{\E}{\mathbb{E}}
\newcommand{\eps}{\varepsilon}
\newcommand{\dd}{\,\mathrm{d}}
\begin{document}
\section*{UnsolvedMath ID: 700617}
\textbf{Source problem code:} catalogue entry
\subsection{Definitions and mathematical statement}
Let \(p\) be a prime. A \(p\)-group here means a group \(G\), possibly infinite, in which every element has finite order equal to a power of \(p\). Such a group has finite exponent if there exists an integer \(e\geq 0\) such that \(g^{p^e}=1\) for every \(g\in G\). The bound may depend on \(G\).

For this draft, interpret \emph{binary finite} as the condition that \(\langle x,y\rangle\) is finite for every \(x,y\in G\). The supplied record does not define this term, so this interpretation requires confirmation against the source. A group \(G\) is \emph{locally finite} if \(\langle S\rangle\) is finite for every finite subset \(S\subseteq G\).

Subject to this terminology qualification, the questions are:
\begin{enumerate}
\item[(a)] Is every binary finite \(2\)-group of finite exponent locally finite?
\item[(b)] For every prime \(p>2\), is every binary finite \(p\)-group of finite exponent locally finite?
\end{enumerate}
Equivalently, for each prime specified in the respective part, does every \(p\)-group \(G\) satisfy the implication
\[
\left[\exists e\in\Z_{\geq 0}\ \forall g\in G,\ g^{p^e}=1\right]
\ \land\ 
\left[\forall x,y\in G,\ |\langle x,y\rangle|<\infty\right]
\quad\Longrightarrow\quad
\left[\forall S\subseteq G,\ |S|<\infty\ \Longrightarrow\ |\langle S\rangle|<\infty\right]?
\]
\subsection{Short English statement}
V. P. Shunkov asks whether a group of bounded prime-power exponent must be locally finite when every two elements generate a finite subgroup, first for the prime \(2\), then for every odd prime. This formulation uses an explicitly provisional interpretation of ``binary finite.'' The supplied registry labels the problem open; that label has not been independently verified.
\subsection{Sources}
\begin{itemize}
\item Local UnsolvedMath registry, supplied complete source record: permanent ID \(700617\), exact source code \texttt{catalogue entry}, title \emph{Kourovka Notebook Problem 9.84}. The supplied status is \texttt{open}.
\item The supplied background attributes Problem 9.84 to V. P. Shunkov, proposed in 1984, and cites \emph{The Kourovka Notebook (Unsolved Problems in Group Theory)}, 9th issue (1984), together with the 21st issue edition (2026), arXiv:1401.0300v48. These bibliographic details are supplied attribution and have not been independently checked.
\item The supplied collection description identifies E. I. Khukhro and V. D. Mazurov as editors of the Kourovka Notebook collection covering Issues 1--21 (1965--2026).
\end{itemize}
\end{document}
